% !TEX root = ../main.tex
\pagebreak
\lecture{3}{2026-09-28}{Woche 3 Relation \& Funktion}
\subsection{Relation}
\begin{definition}
    [\textit{binäre} Relation]\footnote{weil zwischen ZWEI Mengen} die Zuweisung von Elementen einer Menge $\mathbb{D}$ zu einer anderen Menge Z, wobei \textbf{keine} Stetigkeit oder Eindeutigkeit notwendig ist.
    \begin{equation}
        \mathbb{A} = \left\{(x,y) \in \mathbb{D} \times Z \mid \textit{bedingung}\right\} \quad\text{oder Tupel}\quad (\mathbb{D}, Z, \textbf{A})
    \end{equation}
    Aus karthesischem Produkt werden Tupel ausgewählt, die die Bedingung erfüllen. Relationen sind Teilmengen $A \in \mathbb{D} \times Z$
\end{definition}

\begin{definition}
    [Umkehrrelation] kehrung der Mengen und der Bedingung. Sodass:
    \begin{equation}
        R^{-1} \text{ von } (D,Z,R) = (Z,D,R^{-1}) \quad \text{oder} \quad R^{-1} = \{(y,x)  \in A_1 \times A_4 \mid y > x\}
    \end{equation}
\end{definition}

\subsection{Funktion und Eigenschaften}
\begin{definition}
    [Funktion] Relationen, worin elementen einer Menge \textbf{D}  einem einzelnen Element einer anderen Menge \textbf{F} zugewiesen werden. $x \in D$ genau ein $y \in Z$. \textit{D} ist die \textbf{Definitionsmenge} und \textit{Z} ist \textbf{Zielmenge}
    \begin{equation}
        f: \textit{D} \rightarrow Z 
    \end{equation}
\end{definition}

\begin{definition}
    [Graph einer Funktion] Menge aller Punkten, die die Funktionsvorschrift erfüllen.
\end{definition}

Graph von Funktionen mit verschiedener D und Z können gleich aussehen. Die Funktionsvorschrift ist also entscheidend für die Unterscheidung. Z.B. eine Funktion dessen 

\begin{definition}
    [Surjetivität] wenn jedes Element der Zielmenge von elementen der Definitionsmenge erreicht wird. Kein Element der Z-Menge wird übersprungen 
    \begin{equation}
        Z = f(D)
    \end{equation}
\end{definition}

\begin{definition}
    [Interjektivität] wenn alle y aus Z einem einzigen x aus D zugewiesen werden und umgekehrt. Also für $x,x' \in D$
    \begin{equation}
        x \neq x' \implies f(x) \neq f(x')
    \end{equation} 
\end{definition}

\begin{definition}
    [Umkehrfunktion] eine Funktion, deren D und Z Mengen, sowie die Bedingung zu einer gegenteiligen Funktion umgekehrt werden können. Gilt nur für \textit{bijektive} Funktionen.
\end{definition}
Die Eigenschaften \textit{\textbf{Surjektivität}} und \textit{\textbf{Interjektivität}} sind von einander unabhängig. Wenn beide zutreffen ist eine Funktion \textit{\textbf{Bijektiv}}

\begin{definition}
    [Komposition] oder Verschachtelung von Funktionen, wobei die Definitionsmenge einer Funktion $f$ gleich der Zielmenge oder Wertebereich einer anderen Funktion $q$ ist. Schreibweise:
    \begin{equation}
        (f \circ q) = f(q(x)) 
    \end{equation}
\end{definition}

\begin{definition}
    [Bild] einer Menge ist die Menge aller Elementen der Zielmenge, die bei einer Inputmenge A ausgegeben werden.
\end{definition}

\begin{theorem}
    [Komposition und Umkehrfunktion] die Umkehrfunktion einer Komposition ist gleich der getauschten Reihenfolge ihrer Umkehrfunktionen.
    \begin{equation}
        f: D_2 \rightarrow Z_2 \quad \text{und} \quad g: D_1 \rightarrow Z_1 \text{ bijektiv, mit } Z_1 = D_2, \text{ dann gilt }
        (f \circ g)^{-1}(z) = (g^{-1} \circ f^{-1})(z)
    \end{equation}
\end{theorem}